%% ACM Web Conference (WWW) format -- acmart, sigconf. Compile on Overleaf (acmart is preinstalled there).
%% Do NOT change margins, fonts or \vspace (template rules). Course target: 6-8 pages.
%% Every table, figure and number comes from  python scripts/report.py  (tables/, figures/, generated/).
%% Red \TODO markers = things to fill in after the real runs (results prose) or that only you can supply.
\documentclass[sigconf]{acmart}

\usepackage{booktabs}
\newcommand{\TODO}[1]{\textcolor{red}{\textbf{[TODO: #1]}}}
\newcommand{\NB}{\textsc{NeuralBandit}}
\newcommand{\fig}[1]{\IfFileExists{#1.pdf}{\includegraphics[width=\linewidth]{#1}}{\fbox{\TODO{run scripts/report.py}}}}
\usepackage{adjustbox}
%% tables are shrunk to the column/page width if too wide, never enlarged
\newcommand{\tab}[1]{\IfFileExists{#1}{\begin{adjustbox}{max width=\linewidth}\input{#1}\end{adjustbox}}{\fbox{\TODO{run scripts/report.py}}}}
\newcommand{\RefAcc}{\TODO{rho}}
\IfFileExists{generated/macros.tex}{\let\RefAcc\relax% generated by scripts/report.py
\newcommand{\RefAcc}{0.860}
\newcommand{\StatBanditron}{59.6}
\newcommand{\StatLinucb}{71.9}
\newcommand{\StatLints}{71.9}
\newcommand{\StatNbthree}{73.0}
\newcommand{\StatNbtwo}{75.2}
\newcommand{\RegNbtwo}{510.6}
\newcommand{\RegNbtwoTs}{542.2}
\newcommand{\RegNbtwoExpthrees}{490.8}
\newcommand{\RegNbtwoTsExpthrees}{523.7}
\newcommand{\RegNbthree}{559.3}
\newcommand{\RegNbthreeTs}{556.4}
\newcommand{\RegNbthreeExpthrees}{576.3}
\newcommand{\RegNbthreeTsExpthrees}{572.9}
}{}

%% ---- Rights / conference block: acmart sample values; replace with what your instructor gives you ----
\setcopyright{acmlicensed}
\copyrightyear{2026}
\acmYear{2026}
\acmDOI{XXXXXXX.XXXXXXX}
\acmConference[Conference acronym 'XX]{Make sure to enter the correct conference title}{2026}{City, Country}
\acmISBN{978-1-4503-XXXX-X/2026/06}

\begin{document}

\title{Uncertainty-Aware Exploration and Shifting-Expert Selection for NeuralBandit under Concept Drift}
%% Author block: plain text only.
\author{Your Full Name}
\affiliation{\institution{Your University}\city{Your City}\country{Your Country}}
\email{your.email@university.edu}
\renewcommand{\shortauthors}{Surname}

\begin{abstract}
\NB{} (Allesiardo et al., 2014) learns a contextual bandit policy with one neural network per arm and two committee
variants (\NB{}2, \NB{}3) that run networks with different hidden sizes and learning rates and select among them with
EXP3. We reproduce \NB{}2 and \NB{}3 against Banditron, LinUCB and contextual Thompson sampling on Forest Cover Type,
with and without the authors' cyclic concept drift, over 10 seeds. The paper leaves the feedback of \NB{}3's per-arm
selectors unspecified; we show that the natural reading (only the played arm's selector learns) collapses below random
play, and identify the reading of the pseudo-code that reproduces the reported behaviour. We then ask whether
uncertainty-aware exploration and shifting-expert selection improve adaptation to drift, testing Thompson-style arm
selection by Monte-Carlo dropout (H1) and EXP3.S model selection (H2) in a $2\times2$ factorial design, and report
cumulative regret, post-drift recovery time and runtime on Covertype and two further UCI datasets.
Our reproduction matches the original's stationary accuracies to within 1.1 points for NB2, NB3, LinUCB and CTS, and
under drift both committees keep at least 36\% lower cumulative regret than LinUCB. MC-dropout exploration did not
help (it raised \NB{}2's regret by 6\% and its recovery time by 51\%), whereas EXP3.S lowered \NB{}2's regret by
4\% and shortened \NB{}3's post-drift recovery by 29\%. On Pendigits and Shuttle, EXP3.S lowered regret for both
committees (by up to 39\%), but Banditron was best on both, and on the nearly separable Pendigits the committees
failed to recover after drifts, which we trace to saturating sigmoid outputs under the quadratic loss.
Code, configurations and a script that generates every table and figure are released.
\end{abstract}

%% Regenerate at https://dl.acm.org/ccs/ccs.cfm and paste the XML here.
\begin{CCSXML}
<ccs2012>
 <concept><concept_id>10010147.10010257.10010282.10010284</concept_id>
  <concept_desc>Computing methodologies~Online learning settings</concept_desc><concept_significance>500</concept_significance></concept>
 <concept><concept_id>10010147.10010257.10010293.10010294</concept_id>
  <concept_desc>Computing methodologies~Neural networks</concept_desc><concept_significance>300</concept_significance></concept>
</ccs2012>
\end{CCSXML}
\ccsdesc[500]{Computing methodologies~Online learning settings}
\ccsdesc[300]{Computing methodologies~Neural networks}
\keywords{contextual bandits, neural networks, Thompson sampling, MC dropout, EXP3.S, concept drift, reproducibility}

\maketitle

\section{Introduction}
In online advertising and recommendation a system observes a context, chooses one action and sees the reward of that
action only, while user behaviour changes over time. \NB{} \cite{allesiardo2014neuralbandit} trains one multilayer
perceptron per action by importance-weighted gradient descent and, because the right network size and learning rate
are unknown in advance, runs a committee of differently configured networks and selects among them online with EXP3
\cite{auer2002nonstochastic}: one selector for the whole policy (\NB{}2) or one per action (\NB{}3). On Forest Cover
Type the committees matched or beat Banditron, LinUCB and contextual Thompson sampling (CTS), with and without drift.

Two design choices limit adaptation. Exploration is $\epsilon$-greedy with a constant rate $\gamma$ that ignores how
uncertain each network is, and EXP3 competes with the best \emph{fixed} model although after a drift the best model
can change; the original curves show committees needing up to 350{,}000 of the 500{,}000 rounds between drifts to
stabilise.

\paragraph{Research question.} Do uncertainty-aware exploration and shifting-expert selection improve \NB{}'s
adaptation to concept drift? We hypothesise that Thompson-style arm selection by Monte-Carlo (MC) dropout lowers
cumulative regret (H1) and that EXP3.S model selection shortens post-drift recovery (H2).

\paragraph{Contributions.}
(1) A from-scratch reproduction of \NB{}2, \NB{}3 and the original baselines under the original protocol, over 10 seeds
with confidence intervals, documenting every under-specified choice (Sec.~\ref{sec:repro}, Appendix~\ref{app:choices}).
(2) A reproducibility finding: for \NB{}3 the reading in which only the played arm's selector is updated is
self-reinforcing and fails; updating every arm's selector with the round's reward, as Algorithm~3 states, works
(Sec.~\ref{sec:nb3}).
(3) Two modifications, MC-dropout Thompson sampling and EXP3.S, evaluated in a $2\times2$ factorial design on both
committees, with a diagnostic that separates slow base networks from slow model selection, ablations, and additional
datasets (Secs.~\ref{sec:method}--\ref{sec:results}).

\section{Background and Related Work}
\paragraph{Contextual bandits.} At round $t$ a context $x_t$ is revealed, the learner plays $k_t\in[K]$ and observes
only $y_{t,k_t}$. LinUCB \cite{li2010linucb} and CTS \cite{agrawal2013cts} assume rewards linear in $x_t$; Banditron
\cite{kakade2008banditron} is a multiclass perceptron trained from bandit feedback.

\paragraph{\NB{}.} Each arm $k$ has a one-hidden-layer sigmoid network $N^k$ trained with the quadratic loss. With
$\hat k_t=\arg\max_k N^k(x_t)$ the play distribution is $P_t(k)=(1-\gamma)\mathbb 1[k=\hat k_t]+\gamma/K$, and the
played arm's network takes the backpropagation step scaled by $1/P_t(k_t)$ (original Eq.~1), whose expectation is the
full-information step. \NB{}2 runs $M$ such policies (``models'') and an EXP3 over them, with probabilities
$p_i=(1-\gamma_m)w_i/\sum_j w_j+\gamma_m/M$ and reward $y_t$ for the drawn model; every model's network for the
played arm is updated. \NB{}3 keeps one EXP3 per arm: each arm draws a model that scores it, and the
$\epsilon$-greedy rule is applied to these scores.

\paragraph{Neural exploration.} NeuralUCB \cite{zhou2020neuralucb} and NeuralTS \cite{zhang2021neuralts} derive
exploration from gradient features with regret guarantees; Riquelme et al.\ \cite{riquelme2018showdown} compare
approximate posteriors for deep Thompson sampling, including dropout; bootstrapped ensembles
\cite{osband2016bootstrapped} are another approximation. MC dropout is the cheapest of these for a committee of 15
models $\times$ $K$ networks, which is why we test it.

\paragraph{Non-stationarity.} Discounted and sliding-window UCB \cite{garivier2011ucbswitching} and variation-budget
analyses \cite{besbes2014nonstationary} address changing reward distributions; EXP3.S \cite{auer2002nonstochastic}
competes with a \emph{shifting} sequence of experts, the situation a committee faces after a drift.

\section{Reproducing \NB{}}\label{sec:repro}
\paragraph{Protocol.} As in the original: Forest Cover Type (581{,}012 rows, 7 classes = arms); each of the 10
continuous variables is recoded into five equal-frequency binary variables, and the 4 wilderness and 40 soil indicators
are kept (94 binary features, 12 active per row). Reward 1 for the correct class. The shuffled data are replayed in a
loop from a random start; drift relabels the classes by a circular shift every 500{,}000 rounds (shifts accumulate).
$\gamma=0.005$, $\gamma_m=0.1$, and 15 models: hidden sizes $\{1,5,25,50,100\}\times$ learning rates
$\{0.01,0.1,1\}$. We run $T=2.5$M rounds (5 drift segments), 10 seeds (the original averages 10 runs), and measure
regret against an offline MLP trained on the same binary features, $\rho=\RefAcc$ (original: a 93\% classifier).
Banditron's $\gamma$, LinUCB's $\alpha$ and CTS's $v$ are not given in the original; we tune each on two seeds disjoint
from the evaluation seeds with an 8-value grid (Appendix~\ref{app:choices}).

\paragraph{Implementation.} Everything is reimplemented in NumPy/Numba; each context is processed as its list of
active features, so a 2.5M-round committee run takes minutes on a laptop core. All methods see the same stream and
drift schedule for a given seed, so comparisons are paired.

\subsection{An under-specified step in \NB{}3}\label{sec:nb3}
The original does not state what feedback \NB{}3's per-arm selectors receive. If only the played arm's selector is
updated with $y_t$ (the reading of our first implementation), a model that \emph{over}-predicts its arm makes that arm
win the argmax more often and so collects more reward; the arm's selector concentrates on over-predicting models and
the committee ends up playing one arm almost always. Algorithm~3 instead updates \emph{each} arm's EXP3 with the
round's reward for the model it drew, so that every selector is credited for the joint decision. Table~\ref{tab:abl}
(row ``played-arm EXP3 only'') quantifies the difference: with only the played arm's selector updated, \NB{}3's
accuracy over the last 100k rounds falls to 3.7\%, below uniform random play (14.3\%), and its final regret triples
(1{,}749k vs.\ 550k on the same 5 seeds).

\paragraph{Results.} Table~\ref{tab:repro} and Fig.~\ref{fig:repro} reproduce Fig.~1 of the original.
Without drift, accuracies over the last 100{,}000 rounds are within 1.1 points of the original for LinUCB, CTS,
\NB{}3 and \NB{}2, and Banditron is 2.6 points above it; the ordering \NB{}2 $>$ \NB{}3 $>$ LinUCB $\approx$ CTS
$>$ Banditron is the original's. \NB{}2 also has the lowest stationary regret (347k vs.\ 360k for LinUCB, paired
Wilcoxon $p=0.010$). Under drift (Fig.~\ref{fig:repro}, right; Table~\ref{tab:drift}) the linear models lose their
advantage at the first shift: LinUCB and CTS fit one stationary ridge model per arm, take on average 374k and 387k
rounds to recover (5 and 19 of 40 drifts not recovered before the next one) and end with $1.76\times$ and
$2.13\times$ \NB{}2's regret; Banditron, which learns slowly but forgets quickly, ends below both. \NB{}2 ends lowest (511k, 43\% below LinUCB, $p=0.002$) and \NB{}3 second (559k; $+9.5\%$ vs.\ \NB{}2,
95\% CI of the paired difference [13k, 78k], $p=0.049$). Our committees recover in 186k (\NB{}2) and 240k
(\NB{}3) rounds on average, below the up to 350k rounds visible in the original's curves.

\begin{table}[t]
\caption{Covertype with a class shift every 500{,}000 rounds: final cumulative regret ($\times10^3$, mean [95\% CI]),
accuracy over the last 100k rounds, mean recovery time (censored/total drifts), runtime, paired Wilcoxon $p$ against
\NB{}2 (10 seeds).}\label{tab:drift}
\tab{tables/tab_drift.tex}
\end{table}

\begin{table}[t]
\caption{Stationary reproduction: accuracy over the last 100{,}000 rounds, original vs.\ ours (mean [95\% bootstrap
CI] over 10 seeds).}\label{tab:repro}
\tab{tables/tab_repro.tex}
\end{table}

\begin{figure*}[t]
\fig{figures/fig1_reproduction}
\caption{Cumulative regret on Covertype, stationary (left) and with a class shift every 500{,}000 rounds (right, dotted
lines), mean and 95\% band over 10 seeds. Counterpart of Fig.~1 of the original.}\label{fig:repro}
\Description{Two line charts of cumulative regret against the round for Banditron, LinUCB, CTS, NB3 and NB2.}
\end{figure*}

\section{Gaps and Proposed Method}\label{sec:method}
\textbf{G1} Exploration is a constant $\epsilon$-greedy rate that ignores model uncertainty.
\textbf{G2} EXP3 competes with the best \emph{fixed} model; after a drift the best model can change, and the original
committees take up to 350{,}000 rounds to stabilise.
\textbf{G3} One dataset, one synthetic drift protocol, no confidence intervals.
\textbf{G4} Slow recovery may come from the base networks rather than from model selection, which no selector can fix.

\subsection{H1: MC-dropout Thompson sampling}
For networks with at least $h_{\min}=25$ hidden units, each decision draws a dropout mask on the hidden units (rate
$p=0.1$, inverted scaling); the sampled outputs act as one approximate posterior sample per arm, the arm with the
highest sample is played, and with probability $\gamma$ a uniform arm is played instead. Training also uses dropout.
Eq.~1 needs the play probability $P_t(a)$, which has no closed form here; we estimate it from $S=4$ further dropout
passes, $\hat P_t(a)=(1-\gamma)\hat p_t(a)+\gamma/K$. Hidden activations are computed once per round and only the
output layer is resampled, so the extra passes are cheap.

\subsection{H2: EXP3.S model selection}
EXP3.S \cite{auer2002nonstochastic} adds a share of the total weight to every model after each update,
\begin{equation}\label{eq:exp3s}
 w_i\leftarrow w_i\,\exp\!\Big(\frac{\gamma_m\hat x_i}{M}\Big)+\frac{\mathrm{e}\,\alpha}{M}\sum_j w_j,
 \qquad \hat x_i=\mathbb 1[i=m_t]\,\frac{y_t}{p_i},
\end{equation}
so a model that became good after a drift can regain probability quickly; $\alpha=0$ is EXP3. We fix $\alpha=1/T$
before the experiments (the value recommended for a horizon $T$); other values are ablations. In \NB{}3 the same
update is applied to every per-arm selector.

\section{Experimental Setup}
\paragraph{Design.} $2\times2$ factorial: exploration ($\epsilon$-greedy vs.\ MC-dropout TS) $\times$ selector (EXP3
vs.\ EXP3.S), on \NB{}2 and on \NB{}3, under the original drift protocol, 10 seeds, against Banditron, LinUCB and CTS.
\paragraph{Additional datasets.} Pendigits (10{,}992 rows, 16 attributes, $K=10$) and Shuttle (58{,}000 rows,
9 attributes, $K=7$, highly imbalanced) from UCI/OpenML, binarised like Covertype, with
$T=1$M rounds and a class shift every 200{,}000 rounds; \NB{} settings unchanged (not re-tuned), baselines tuned per
dataset as above.
\paragraph{Metrics.} Cumulative regret $R_T=\rho T-\sum_t y_{t,k_t}$; accuracy over the last 100{,}000 rounds;
\emph{recovery time}: rounds after a drift until the 20{,}000-round moving accuracy reaches 95\% of its level over the
50{,}000 rounds before the drift (a drift not recovered before the next one counts as the full period, and the number
of such censored drifts is reported); runtime per run.
\paragraph{Statistics.} Paired over seeds: mean difference, 95\% bootstrap CI, exact Wilcoxon signed-rank test.
A difference is called lower/higher only if its CI excludes zero; comparisons are not corrected for multiplicity.
\paragraph{Diagnostic (G4).} In simulation the true label is known, so we also log how often each model's own greedy
arm is correct and plot the best single model in hindsight per drift segment against the realised policy: if the
best model itself recovers slowly, selection cannot be the bottleneck.

\section{Results}\label{sec:results}
\begin{table*}[t]
\caption{Covertype with drift: \NB{}2 extensions ($2\times2$) and LinUCB. Final cumulative regret ($\times10^3$, mean
[95\% CI]), accuracy over the last 100k rounds, mean recovery time (censored/total drifts), runtime, paired Wilcoxon
$p$ against \NB{}2 (10 seeds).}\label{tab:ext2}
\tab{tables/tab_ext_nb2.tex}
\end{table*}
\begin{table*}[t]
\caption{As Table~\ref{tab:ext2} for \NB{}3.}\label{tab:ext3}
\tab{tables/tab_ext_nb3.tex}
\end{table*}
\begin{figure}[t]
\fig{figures/fig_ext_nb2}
\caption{Cumulative regret of the \NB{}2 variants and LinUCB under drift (mean, 95\% band, 10 seeds).}\label{fig:ext}
\Description{Line chart of cumulative regret against the round for NB2, its three variants and LinUCB.}
\end{figure}
\begin{figure}[t]
\fig{figures/fig_diag_nb2}
\caption{Diagnostic: moving accuracy of the realised \NB{}2 policy and of the best single model in hindsight per
drift segment (simulation only; uses true labels).}\label{fig:diag}
\Description{Two accuracy curves against the round, dropping at every drift.}
\end{figure}

\paragraph{H1 (MC-dropout TS): not supported.} On \NB{}2, MC-dropout Thompson sampling raised final regret by 31.6k
($+6.2\%$, 95\% CI [3.5k, 53.8k], $p=0.064$), raised mean recovery time by 95.6k rounds ($+51\%$, CI [64.8k, 129.9k],
$p=0.002$) and lowered accuracy over the last 100k rounds from 68.3\% to 66.2\% (Table~\ref{tab:ext2}). On \NB{}3
it made no clear difference ($-0.5\%$, CI [$-42.4$k, 38.5k], $p=0.85$; Table~\ref{tab:ext3}). After a drift the
networks are confidently wrong rather than uncertain, so dropout noise does not point exploration at the arms that
changed; it adds noise to every decision, and its estimated propensities and dropout training add variance to the
Eq.~1 updates. We read this as a limit of dropout as a posterior under non-stationarity, not as evidence against
uncertainty-aware exploration in general.
\paragraph{H2 (EXP3.S): partly supported.} EXP3.S lowered \NB{}2's final regret by 19.8k ($-3.9\%$, CI [$-41.3$k,
$-4.9$k], $p=0.020$), the lowest regret of all methods, but did not change its recovery time ($+14.3$k rounds, CI
[$-17.3$k, 45.6k], $p=0.49$). On \NB{}3 the effect was the reverse: mean recovery time fell by 69.9k rounds
($-29\%$, CI [$-99.8$k, $-36.6$k], $p=0.010$) and unrecovered drifts from 6 to 1 of 40, while final regret did not
clearly change ($+3.0\%$, CI [$-0.6$k, 31.6k], $p=0.11$). \NB{}3 must re-identify a good model in each of $K$
selectors after a shift, which is where EXP3.S's weight sharing helps recovery; \NB{}2 has one selector, where the
sharing mostly avoids locking onto a model that later degrades.
\paragraph{Interaction and diagnostic.} The two changes did not interact: the $2\times2$ interaction contrast on
final regret was $+1.3$k (CI [$-25.1$k, 36.1k]) for \NB{}2 and $-0.6$k (CI [$-19.2$k, 18.1k]) for \NB{}3, and
\NB{}3 with both had the shortest recovery of any committee (155k rounds, $-35\%$, $p=0.010$). In the diagnostic
(Fig.~\ref{fig:diag}) the best single model in hindsight averages 69.3\% moving accuracy against 65.4\% for the
realised \NB{}2 policy (69.8\% vs.\ 66.2\% with EXP3.S). Selection therefore costs about 4 points, and EXP3.S
narrows that gap only slightly; the best model itself drops at every shift and climbs back over a few hundred
thousand rounds, so most of the recovery time is the networks relearning, which no selector can remove (G4).
\paragraph{Comparison with LinUCB.} On Covertype, every committee variant ends with lower cumulative regret than
LinUCB under drift ($p=0.002$ for all eight): from 45\% lower for \NB{}2+EXP3.S ($-408.8$k, CI [$-422.0$k, $-395.3$k]) to 36\%
lower for \NB{}3+EXP3.S. The price is runtime: about 9--14 minutes per 2.5M-round run against 0.6 minutes for LinUCB, with MC
dropout adding 0--40\% to a committee's runtime.
\paragraph{Ablations.} Table~\ref{tab:abl} (5 seeds, descriptive). Clipping the importance weight is essential:
the literal unclipped update doubles \NB{}2's regret (1{,}001k vs.\ 502k) and drops its final accuracy to 42.8\%,
whereas clip 10 and clip 100 are similar (505k vs.\ 502k). Dividing by the committee's propensity instead of each
model's own raised regret for both committees (\NB{}2 539k, \NB{}3 640k) although it shortened recovery. Small
EXP3.S shares help and large ones hurt: $\alpha=10^{-6}$ gave 487k and $\alpha=10^{-4}$ 560k. Every MC-dropout
setting we tried ($p\in\{0.05,0.1,0.2\}$, $S\in\{1,4,8\}$) had higher regret than \NB{}2 (533k--551k), so the
negative H1 result does not hinge on these choices.
\paragraph{Additional datasets.} Table~\ref{tab:extra} and Fig.~\ref{fig:extra}. The ordering of the baselines
changes: Banditron has the lowest regret on both datasets (Pendigits 228k, Shuttle 78k), and on Pendigits LinUCB (488k)
also beats every committee (541k--693k), while on Shuttle all committees beat LinUCB (310k) and CTS (466k). On
Pendigits, \NB{}2 learns before the first drift but then stays near chance (14.3\% over the last 100k rounds; 26 of
40 drifts not recovered). Pendigits is almost separable ($\rho=0.98$) and is replayed about 90 times, so the networks
become very confident; with sigmoid outputs and the quadratic loss the gradient $\delta=(y-N)N(1-N)$ vanishes for
confident outputs, so a network that is confidently wrong after a shift barely learns. Banditron, a perceptron
without a squashing output, does not have this problem. EXP3.S again helps: it lowered \NB{}2's regret by 7.8\% on
Pendigits ($p=0.002$) and by 38.5\% on Shuttle (CI [$-111.5$k, $-61.6$k], $p=0.002$), and it nearly doubled
\NB{}3's final accuracy on Pendigits (48.5\% vs.\ 27.1\%). MC dropout lowered \NB{}2's regret on Pendigits by
2.9\% ($p=0.037$) and made no clear difference on Shuttle ($p=0.19$).

\begin{table*}[t]
\caption{Additional datasets with a class shift every 200{,}000 rounds ($T=1$M): final cumulative regret ($\times10^3$,
mean [95\% CI]) and accuracy over the last 100k rounds, 10 seeds. Paired tests are in \texttt{results/REPORT.md}.}
\label{tab:extra}
\tab{tables/tab_extra.tex}
\end{table*}
\begin{figure*}[t]
\begin{minipage}{0.49\linewidth}\fig{figures/fig_pendigits}\end{minipage}\hfill
\begin{minipage}{0.49\linewidth}\fig{figures/fig_shuttle}\end{minipage}
\caption{Cumulative regret on Pendigits (left) and Shuttle (right) with a class shift every 200{,}000 rounds (mean,
95\% band, 10 seeds).}\label{fig:extra}
\Description{Two line charts of cumulative regret for Banditron, LinUCB, CTS, NB3, NB2, NB2+MCD-TS and NB2+EXP3.S.}
\end{figure*}

\begin{table*}[t]
\caption{Ablations on Covertype with drift (seeds 0--4; descriptive, so $p$ values are not interpreted). Runtimes
are omitted because the machine was suspended during some ablation runs.}\label{tab:abl}
\tab{tables/tab_ablations.tex}
\end{table*}

\section{Discussion, Limitations and Future Work}
The reproduction supports the original's central claim: under the cyclic drift protocol, committees of small
networks adapt far better than linear contextual bandits, whose single ridge model accumulates stale evidence. The additional datasets qualify this: on Shuttle the committees still beat the linear models, but on
the nearly separable Pendigits they fail to recover from drift, because confident sigmoid outputs stop learning under
the quadratic loss, and a simple perceptron is best on both. Of our two modifications, the cheap one (EXP3.S, one
extra line in the selector update) helped on all three datasets and the uncertainty-based one did not. The diagnostic explains why selection-side changes give modest gains: the base networks, not the selector,
set the pace of recovery. The \NB{}3 finding is a reminder that a committee's credit-assignment rule can decide
whether it works at all, and that pseudo-code steps such as ``update each arm's EXP3'' need to be stated exactly.
\paragraph{Limitations.} (i) The drift is a synthetic class permutation: it changes the label--arm mapping, not the
covariate distribution. (ii) Additional datasets are small and replayed many times. (iii) Choices the original leaves
open (Appendix~\ref{app:choices}) affect results. (iv) Dropout is not a calibrated posterior and its play probability
is estimated. (v) NeuralUCB and NeuralTS were not run (laptop budget). (vi) Comparisons are not corrected for
multiplicity. (vii) Recovery time is measured against each method's own pre-drift level, so a method that recovers
quickly to a low level (Banditron, 78k rounds) can look good on this measure alone. (viii) We kept the original
quadratic loss; its saturation is our explanation for the Pendigits failure, and a cross-entropy loss was not tested.
\paragraph{Future work.} A cross-entropy loss and change
detection for the committees' networks; natural covariate drift and real logged-bandit data; NeuralUCB/NeuralTS baselines; learned or LLM-based context encoders for text-based bandits.

\section{Conclusion}
We reproduced \NB{}2 and \NB{}3 on Covertype to within about one accuracy point of the original and confirmed
that on Covertype they beat Banditron, LinUCB and CTS under concept drift, after identifying the reading of \NB{}3's selector
update that makes it work. EXP3.S model selection lowered cumulative regret on all three datasets and shortened \NB{}3's
post-drift recovery on Covertype, whereas MC-dropout Thompson sampling did not reliably improve either committee; on
nearly separable data the committees themselves stop adapting. Faster adaptation will
require making the base networks relearn faster after a change, for example with a
cross-entropy loss, or by detecting changes and resetting or re-weighting part of the committee.

\section*{Ethics Statement}
Adaptive recommendation and advertising systems can reinforce engagement-driven feedback loops, and exploration
exposes users to lower-reward items. Our experiments use public, non-personal benchmark data and a synthetic drift
process; methods that adapt faster to changing behaviour should be evaluated for fairness and user welfare before
deployment on personal data.

\begin{acks}
Final project for the course \emph{Deep Learning for Data Science}.
\end{acks}

\bibliographystyle{ACM-Reference-Format}
\bibliography{refs}

\appendix
\section{Choices the Original Leaves Open}\label{app:choices}
\textbf{Networks.} Biases included; weights $\sim U(-0.5,0.5)$; sigmoid output, quadratic loss,
$\delta=(y-N)N(1-N)$. \textbf{Eq.~1 propensity.} Each model is an \NB{}1 instance and divides by its own
$P_t(k_t)$ (\NB{}2 and \NB{}3); the committee's propensity is an ablation. \textbf{Clipping.} The weight
$1/P_t$ (up to $K/\gamma=1400$) is clipped at 100; no clip and clip 10 are ablations. \textbf{\NB{}3 selectors.} Every
arm's EXP3 is updated with $y_t$ for the model it drew (Sec.~\ref{sec:nb3}). \textbf{Drift.} Arm of the correct answer
$=(\text{label}+\lfloor t/500{,}000\rfloor)\bmod K$. \textbf{Baselines.} LinUCB and CTS use disjoint ridge models with
$A=I+\sum xx^\top$ and the same binary features plus a bias; CTS samples $\tilde\theta^\top x\sim
\mathcal N(\hat\theta^\top x, v^2x^\top A^{-1}x)$. \textbf{Tuning} (Covertype, seeds 100--101, 1M rounds, one drift; 8
log-spaced values each, highest mean reward selected): Banditron $\gamma\in[0.001,0.5]$, selected 0.02; LinUCB
$\alpha\in[0.03,100]$, selected 3; CTS $v\in[0.003,10]$, selected 1. No selected value lies at the edge of its grid;
the full grid is in \texttt{paper/tables/tab\_tuning\_covtype.tex}.
\textbf{Stream.} Shuffled per seed, random start, replayed in a loop (about 4 passes at $T=2.5$M).


\end{document}
